% 附录 E PYDB Perl API
\biappendix{PYDB Perl API}{PYDB Perl API}
\label{app:e}

本附录介绍如何使用 IC Validator 工具提供的 Perl 应用程序编程接口（API）访问 IC Validator 错误数据库（PYDB）。

PYDB Perl API 使你能直接访问 IC Validator 工具产生的错误数据。例如，你可以：
\begin{itemize}
  \item 为特定违例生成错误信息；
  \item 统计违例数量以供跟踪。
\end{itemize}

数据库与 Perl API 包括以下各节：
\begin{itemize}
  \item PYDB Perl API 概述（PYDB Perl API Overview）
  \item 访问 PYDB（Accessing the PYDB）
  \item 常见任务示例（Examples of Common Tasks）
  \item PYDB 模式（PYDB Schema）
\end{itemize}

% ============================================================
\section{PYDB Perl API 概述}
\subsection*{PYDB Perl API Overview}

使用 PYDB Perl API，你可以访问 IC Validator 工具发现的违例信息，并按需处理数据。IC Validator 仅将含错误数据的违例与单元存入 PYDB。\cmd{LAYOUT\_ERRORS} 文件通常显示 PYDB 中可用的大部分信息。

IC Validator 产生的错误数据以表的形式存储在 SQLite 数据库中；该数据库即为 PYDB。表的结构及其相互关系称为数据库模式（schema）。模式信息见「PYDB 模式」。

要正确查询复杂数据库，必须了解数据库模式。可用 SQLite 命令访问数据库并返回信息；这些命令可能较复杂，取决于所需返回的信息。

为便于在无需学习数据库模式的情况下访问错误数据，所提供的 PYDB Perl API 可连接并检索 PYDB。与使用 SQLite 命令相比，它使与错误数据库的交互容易得多。其内置功能包括：
\begin{itemize}
  \item 连接到 PYDB 服务器；
  \item 提交查询并确保多个数据库服务器不会相互竞争；
  \item 提供高级函数以检索 PYDB 特定信息，而无需编写复杂的 SQLite 查询；
  \item 处理结果。
\end{itemize}

PYDB Perl API 位于 \cmd{ICV\_INSTALL\_DIRECTORY/contrib/pydb/}。子目录见表~\ref{tab:app-e-pydb-dirs}。

\begin{table}[htbp]
\centering
\caption{pydb 目录的子目录}
\label{tab:app-e-pydb-dirs}
\begin{tabular}{@{}>{\ttfamily\raggedright\arraybackslash}p{0.22\textwidth} p{0.72\textwidth}@{}}
\toprule
\textbf{\rmfamily 子目录} & \textbf{说明} \\
\midrule
bin/ &
含示例脚本。这些脚本展示如何使用 PYDB 模块。 \\
\midrule
lib/ &
含公共库模块。每个文件有简短说明，后接该模块提供的全部类方法的详细说明。 \\
\bottomrule
\end{tabular}
\end{table}

% ============================================================
\section{访问 PYDB}
\subsection*{Accessing the PYDB}

PYDB Perl API 提供高级接口，使你能处理错误数据，并在需要时对错误数据库运行直接 SQLite 查询。

\begin{noteBox}
PYDB 模式与服务器通信协议可能随版本变化。因此，访问 PYDB 时必须使用生成该 PYDB 的同一 IC Validator 版本所附带的 API。
\end{noteBox}

使用 IC Validator 提供的 PYDB Perl API 访问 PYDB，请按下列步骤：

\begin{enumerate}
  \item \textbf{加载 PYDB 库。}\\
在任何数据库命令之前，在 Perl 脚本中包含下列代码。该代码从 IC Validator 安装导入所需模块：
\begin{lstlisting}
push @INC, $ENV{"ICV_INSTALL_DIRECTORY"}."/contrib/pydb/lib";
require "PYDB.pm";
\end{lstlisting}

  \item \textbf{连接数据库。}\\
在 Perl 脚本中包含下列代码以连接数据库：
\begin{lstlisting}
$dbh = PYDB::Connect($db_path,$db_name);
if (! $dbh) {
   print "ERROR: Error connecting to database: $PYDB::errstr\n";
   exit;
}
\end{lstlisting}
若连接成功，返回数据库句柄。

  \item \textbf{执行数据库查询。}\\
使用 PYDB Perl API 可访问错误信息并按需处理数据。仅含错误数据的违例与单元存入数据库。\cmd{LAYOUT\_ERRORS} 文件通常显示 PYDB 中可用的大部分信息。示例见「常见任务示例」。

  \item \textbf{断开数据库连接。}\\
在 Perl 脚本中包含下列代码以断开连接：
\begin{lstlisting}
$dbh->Disconnect();
\end{lstlisting}
\end{enumerate}

% ============================================================
\section{常见任务示例}
\subsection*{Examples of Common Tasks}

下列示例展示一些常见任务。

\subsection{列出数据库中的违例}
\subsubsection*{List violations in the database}

\begin{lstlisting}
# All violations stored in the PYDB
foreach $vio ($dbh->Violations()) {
   print "Violation: \"".$vio->Comment()."\"\n";
}

# Find a specific violation
$vio = $dbh->GetViolationByComment("M1.S1.1: Metal 1 space");
if ($vio) {
   print "Found M1 violation, with ".$vio->ErrorCount().
      " error(s).\n";
}

# Find violations that match a regular expression
foreach $vio ($dbh->GetViolationsByCommentRegex(qr/M1\..*/)) {
   print "M1 rule: \"".$vio->Comment()."\"\n";
}
\end{lstlisting}

\subsection{列出数据库中的单元}
\subsubsection*{List cells in the database}

\begin{lstlisting}
# All cells stored in the PYDB
foreach $cell ($dbh->Cells()) {
   print "Cell: \"".$cell->CellName()."\"\n";
}

# Find a specific cell
$cell = $dbh->GetCellByName("TOP");
if ($cell) {
   print "Found top cell.\n";
}

# Find cells that match a regular expression
foreach $cell ($dbh->GetCellsByNameRegex(qr/STDCELL.*/)) {
   print "Found cell: \"".$cell->CellName()."\"\n";
}
\end{lstlisting}

\subsection{访问特定违例中的错误}
\subsubsection*{Access errors in a specific violation}

\begin{lstlisting}
$vio = $dbh->GetViolationByComment("M1.S1.1: Metal 1 space");
if ($vio) {

    foreach $command ($vio->Commands()) {
       print "Command ".$command->Comment().
          " produced ".$command->NumErrorsUnlimited().
          " errors, details are available for ".
          $command->NumErrors()."\n";

        foreach $cell ($command->Cells()) {
           # Errors are stored by command/cell, the error
           # iterator can be used to access them
           $iter = $dbh->ListErrors($command, $cell);
           if (! $iter) {
              print "ERROR: $PYDB::errstr\n";
              $dbh->Disconnect();
              exit;
           }

            while ($error = $iter->NextError()) {
               # Bounding box of error
               @bbox = $error->GetBBox();
               # Error polygons
               foreach $polygon ($error->Polygons()) {
                  # List of points making up the polygon
                  @points = $polygon->Points();
               }
               # Access to specific data fields by column name
               $distance = $error->Get("Distance");
               # Certain types of errors might have child errors
               foreach $child_error ($error->Children()) {
                  # ...
               }
            }
        }
    }
}
\end{lstlisting}

\subsection{执行直接 SQLite 查询}
\subsubsection*{Perform a direct SQLite query}

Perl API 提供与多数 Perl 安装自带的 DBI 模块类似的接口。

\begin{lstlisting}
# Prepare the query
$sth = $dbh->prepare("SELECT CellID, CellName FROM pyCellNames");
if (! $sth) {
   print "ERROR: $PYDB::errstr\n";
   $dbh->Disconnect();
   exit;
}

# Execute the query
if (! $sth->execute()) {
   print "ERROR: $PYDB::errstr\n";
   $dbh->Disconnect();
   exit;
}

# Retrieve the data
while ($ref = $sth->fetchrow_hashref()) {
   print "Cell name=".$ref->{'CellName'}.
      ", id=".$ref->{'CellID'}."\n";
}
\end{lstlisting}

% ============================================================
\section{PYDB 模式}
\subsection*{PYDB Schema}

在 PYDB 模式中，错误存储在错误表中。
\begin{itemize}
  \item 每个表由可引用若干相关列的数据行组成，类似电子表格。
  \item 每个表至少需要一个唯一键以标识各条记录。某些表引用多个键。
\end{itemize}

每种命令类型至少有一个表；含有子错误的命令还有额外表。错误分类信息也存储在 PYDB 的若干表中。

表~\ref{tab:app-e-pydb-org} 展示 PYDB 中一些常用表。

\begin{longtable}{@{}>{\ttfamily\raggedright\arraybackslash}p{0.32\textwidth} p{0.62\textwidth}@{}}
\caption{PYDB 组织结构}\\
\label{tab:app-e-pydb-org}\\
\toprule
\textbf{\rmfamily 表} & \textbf{存储的信息} \\
\midrule
\endfirsthead
\toprule
\textbf{\rmfamily 表} & \textbf{存储的信息} \\
\midrule
\endhead
\bottomrule
\endfoot

pyCellNames &
含单元信息。每个命令可从不同单元写入错误。 \\
\midrule

pyCommandCells &
列出某命令含错误的单元。 \\
\midrule

pyCommands &
对产生错误的每个 runset 命令含一条记录。多个命令（例如 \cmd{internal1()} 与 \cmd{external2()}）可写入同一违例注释。 \\
\midrule

pyErrors\_\textit{<command type>} &
含各条错误数据。特定命令类型的错误表动态创建，并含该命令类型特有的数据列。 \\
\midrule

pyErrorTableFormat &
定义错误表的数据列。 \\
\midrule

pyErrorTables &
定义与某命令类型关联的错误表层次。 \\
\midrule

pyPolygons\_\textit{<command type>} &
含各条错误形状。 \\
\midrule

pyRunInfo &
含一般信息。 \\
\midrule

pyViolations &
定义 PYDB 中存储错误的违例规则注释。 \\
\end{longtable}
